% Arquivo LaTeX de exemplo de dissertação/tese a ser apresentada à CPG do IME-USP
%
% Criação: Jesús P. Mena-Chalco
% Revisão: Fabio Kon e Paulo Feofiloff
% Adaptação para UTF8, biblatex e outras melhorias: Nelson Lago
%
% Except where otherwise indicated, these files are distributed under
% the MIT Licence. The example text, which includes the tutorial and
% examples as well as the explanatory comments in the source, are
% available under the Creative Commons Attribution International
% Licence, v4.0 (CC-BY 4.0) - https://creativecommons.org/licenses/by/4.0/

% ----------------------------------------------------------
\makeatletter
\def\input@path{{aux/}}
\makeatother

\documentclass[a4paper,12pt,twoside,english,brazilian]{book}

\usepackage{imegoodies}
\usepackage[thesis]{aux/imelooks}
\usepackage{imakeidx}

%\nocolorlinks

% Caminhos para figuras
\graphicspath{{img/}}

% Comandos de idioma
\babeltags{br = brazilian, en = english}

% Bibliografia
\addbibresource{bibliografia.bib}
%\nocite{*} % incluir todas as referências

% Hifenizações específicas
\babelhyphenation{documentclass latexmk soft-ware clsguide}
\babelhyphenation[brazilian]{Fu-la-no}
\babelhyphenation[english]{what-ever}

% Título e autoria
\title{Título do trabalho}[um subtítulo]
\translatedtitle{Title of the document}[a subtitle]
\author[fem]{Nome Completo}

% Orientação
\def\profa{Prof\kern.02em.\kern-.07emª\kern.07em}
\def\dra{Dr\kern-.04em.\kern-.11emª\kern.07em}
\orientador[fem]{\profa{} \dra{} Fulana de Tal}
\coorientador{Prof. Dr. Ciclano de Tal}
\coorientador[fem]{\profa{} \dra{} Beltrana de Tal}

% Banca
\banca{
  \profa{} \dra{} Fulana de Tal (orientadora) -- IME-USP,
  Prof. Dr. Ciclano de Tal -- IME-USP,
  \profa{} \dra{} Convidada de Tal -- IMPA,
}

% Tipo de trabalho e programa
\tipotese{
  mestrado,
  %doutorado,
  %tcc,
  %definitiva,
  %quali,
  programa={Ciência da Computação},
}

% Defesa
\defesa{
  local={São Paulo},
  data=2017-08-10,
}

% Licença
%\direitos{CC-BY-NC-ND}
%\direitos{Autorizo a reprodução e divulgação total ou parcial deste
%          trabalho, por qualquer meio convencional ou eletrônico,
%          para fins de estudo e pesquisa, desde que citada a fonte.}
%\direitos{I authorize the complete or partial reproduction and disclosure
%          of this work by any conventional or electronic means for study
%          and research purposes, provided that the source is acknowledged.}
\direitos{CC-BY}

% Ficha catalográfica (preencher com código gerado)
\fichacatalografica{}

% -------------------------------------------
% DOCUMENTO
% -------------------------------------------

\begin{document}

% Páginas iniciais
\frontmatter
\pagestyle{plain}
\onehalfspacing
\maketitle

% Dedicatória
\begin{dedicatoria}
Esta seção é opcional.
\end{dedicatoria}

\pagenumbering{roman}

% Agradecimentos
\chapter*{Agradecimentos}
\epigrafe{Do. Or do not. There is no try.}{Mestre Yoda}
Texto texto texto...

% Resumo e abstract
\input{conteudo/resumo_abstract}

% Listas manuais
\newcommand\disablenewpage[1]{{\let\clearpage\par\let\cleardoublepage\par #1}}
\bgroup
\raggedbottom

\disablenewpage{\chapter*{Lista de abreviaturas}}
\begin{tabular}{rl}
   CFT & Transformada contínua de Fourier\\
   DFT & Transformada discreta de Fourier\\
  EIIP & Potencial de interação elétron-íon\\
  STFT & Transformada de Fourier de tempo reduzido\\
  ABNT & Associação Brasileira de Normas Técnicas\\
   URL & Uniform Resource Locator\\
   IME & Instituto de Matemática e Estatística\\
   USP & Universidade de São Paulo
\end{tabular}

\disablenewpage{\chapter*{Lista de símbolos}}
\begin{tabular}{rl}
  $\omega$ & Frequência angular\\
  $\psi$ & Função de análise wavelet\\
  $\Psi$ & Transformada de Fourier de $\psi$\\
\end{tabular}

% Listas automáticas
\clearpage
\disablenewpage{\listoffigures}
\disablenewpage{\listoftables}
\disablenewpage{\listof{program}{\programlistname}}

% Sumário
\tableofcontents
\egroup

% Índice remissivo (remissivas iniciais)
\index{Inglês|see{Língua estrangeira}}
\index{Figuras|see{Floats}}
\index{Tabelas|see{Floats}}
\index{Código-fonte|see{Floats}}
\index{Subcaptions|see{Subfiguras}}
\index{Equações|see{Modo matemático}}
\index{Rodapé, notas|see{Notas de rodapé}}
\index{Captions|see{Legendas}}

% Capítulos
\mainmatter
\pagestyle{mainmatter}
\singlespacing

\pagestyle{unnumberedchapter}
\input{conteudo/00_exemplo_introducao}
\pagestyle{mainmatter}
\input{conteudo/01_exemplo_normas_ime}
\input{conteudo/02_exemplo_usando_o_modelo}
\input{conteudo/03_exemplo_latex}
\input{conteudo/04_exemplo_tutorial}
\input{conteudo/05_exemplo_exemplos}

% Apêndices
\cleardoublepage
\pagestyle{appendix}
\appendix
\addappheadtotoc
\input{conteudo/apendice_exemplo_faq}

% Anexos
\cleardoublepage
\annex
\addappheadtotoc
\input{conteudo/anexo_exemplo_imegoodies}
\input{conteudo/anexo_exemplo_pseudocodigo}

% Referências
\backmatter
\pagestyle{backmatter}
\addtocontents{toc}{\vspace{2\baselineskip plus .5\baselineskip minus .5\baselineskip}}

\printbibliography[
  title=\refname\label{sec:bib},
  heading=bibintoc,
]

% Índice remissivo (opcional)
\printindex

\end{document}
